\scanpage{142}
la situation pour $G_1$, $G'_1$, $(P_1,\xi_1)$, $(Q_1,\eta_1)$ etc. [Attention, on perd le fait que $\unsure{S_2|S|S_1|Q_1}$ soit réduit, $S_1$ ne l'étant plus]. $P_1$ s'\uline{identifie} alors $G_1$\inserted{(muni du \unsure{pt distingué|point marqué|marquage} $\varepsilon_1$)}, $Q_1$ à $G'_1\times_{S_1}G'_1$ muni de $(\varepsilon'_1,\varepsilon'_1)$, \deleted{je dis que} enfin $(P'_1,\xi'_1)$ s'identifie aussi à $(G'_1,\varepsilon'_1)$. Il suffit de voir que $(P'_1,\xi'_1)\to(\unsure{Q'_1|Q_1},\eta_1)$ est l'application diagonale $G'_1\unsure{-|\to}G'_1\times_{S_1}G'_1$.

\uline{Or} \unsure{remontons} les définitions, et désignons par $w$ l'isom. $(P,\xi)\times_u G'\longrightarrow(P,\xi)\times_vG'$, $w_1$ l'isom. déduit par extension de la base à $S_1$, notre assertion signifie que $w_1$ transforme la section ``naturelle'' de $P_1\times_{u_1}G'_1$ dans celle de $P_1\times_{v_1}G'_1$, i.e. qu'on a commutativité dans

\diagram{\begin{tikzcd}[ampersand replacement=\&,column sep=small]
P=S_1\arrow[r]\arrow[rr,bend right=16] \& (P\times_uG')\times_S S_1\arrow[r,"{w_1=w\times_{\unsure{S|S,S}} S'}"] \& (P\times_v G')\times_S S_1
\end{tikzcd}}

ou encore qu'on a commutativité de

\diagram{\begin{tikzcd}[ampersand replacement=\&]
P\arrow[r,"\mathrm{can}"]\arrow[rr,bend right=15,"\mathrm{can}"']\& P\times_uG'\arrow[r,"w"]\& P\times_vG'
\end{tikzcd}}

\uline{Or cela résulte} de \uline{l'unicité} de $P\longrightarrow P\times_vG'$ \uline{transformant} \deleted{\unsure{$\xi$ en $\xi'$}} pt marqué en pt marqué \ldots.

% archive sheet number: 141

